%!TEX program = xelatex
%!TEX encoding = UTF-8
%Definições de Ano, Semestre, número e data da avaliação
\def\numeroteste{4}
\def\modulo{2}
\def\dataavaliacao{01/11/2023}

\documentclass[12pt]{exam}

\usepackage[singlelinecheck=false]{caption}
\usepackage{textpos}
\usepackage{amssymb}
\usepackage{amsmath,amsfonts,amsthm,amstext}
\usepackage[justification=centering]{caption}
\usepackage[brazil]{babel}
% \usepackage[latin1]{inputenc}
\usepackage{graphicx}
\graphicspath{{../../figuras_geogebra/}{../../Pictures/}}
\usepackage{enumitem}
\usepackage{multicol}
\usepackage{answers}
\usepackage{qrcode}

\usepackage{pifont}
\newcommand{\cmark}{\ding{228}}%

% requires version 0.3 of the package
\usepackage[customcolors,beamer]{hf-tikz}

\tikzset{style green/.style={
        set fill color=green!50!lime!60,
        set border color=white,
    },
    style cyan/.style={
        set fill color=cyan!90!blue!60,
        set border color=white,
    },
    style orange/.style={
        set fill color=orange!80!red!60,
        set border color=white,
    },
    hor/.style={
        above left offset={-0.15,0.31},
        below right offset={0.15,-0.125},
        #1
    },
    ver/.style={
        above left offset={-0.1,0.3},
        below right offset={0.15,-0.15},
        #1
    }
}

\usepackage{tikz,ifthen,tikz-3dplot}
\usepackage{siunitx}
\usetikzlibrary{circuits.ee.IEC}
\usetikzlibrary{lindenmayersystems}
\usetikzlibrary[shadings]
\usepackage{pgfplots}
\pgfplotsset{compat=1.15}
\usepackage{mathrsfs}
\usetikzlibrary{arrows}
\usepackage{tcolorbox}
\newcounter{saveenumi}
\newcommand{\seti}{\setcounter{saveenumi}{\value{enumi}}}
\newcommand{\conti}{\setcounter{enumi}{\value{saveenumi}}}

\resetcounteronoverlays{saveenumi}


\renewcommand{\qedsymbol}{$\blacksquare$}

%\setbeamertemplate{itemize items}[ball]

%\setbeamertemplate{itemize items}[default]
%-------------amatrix
% Augmented matrix.  Usage (note the argument does not count the aug col):
% \begin{amatrix}{2}
    %   1  2  3 \\  4  5  6
    % \end{amatrix}
\newenvironment{amatrix}[1]{%
    \left[\begin{array}{@{}*{#1}{r}|r@{}}
    }{%
    \end{array}\right]
}

\def\ano{2023}
\def\semestre{2}
\def\turma{11}
\newcommand{\nsub}{\varsubsetneq}
\newcommand{\vaz}{\emptyset}
\newcommand{\dom}{{\rm dom\,}}
\newcommand{\im}{{\rm Im\,}}
\newcommand{\sub}{\subseteq}
\newcommand{\n}{\mathbb{N}}
\newcommand{\z}{\mathbb{Z}}
\newcommand{\rac}{\mathbb{Q}}
\newcommand{\real}{\mathbb{R}}
\newcommand{\complex}{\mathbb{C}}
\newcommand{\cp}[1]{\mathbb{#1}}
\newcommand{\ch}{\mbox{\textrm{car\,}}\nobreak}
\newcommand{\p}[1]{\mbox{\textrm{posto(#1)}}}
\newcommand{\vesp}[1]{\vspace{ #1  cm}}
\newcommand{\compcent}[1]{\vcenter{\hbox{$#1\circ$}}}
\newcommand{\comp}{\mathbin{\mathchoice
{\compcent\scriptstyle}{\compcent\scriptstyle}
{\compcent\scriptscriptstyle}{\compcent\scriptscriptstyle}}}
\newcommand\T{\rule{0pt}{2.6ex}}
\newtheorem{definicao}{Defini\c{c}\~ao}[section]
\newtheorem{definicoes}{Defini\c{c}\~oes}[section]
\newtheorem{exemplo}{Exemplo}[section]
\newtheorem{exemplos}{Exemplos}[section]
\newtheorem{exercicio}{Exerc{\'\i}cio}
\newtheorem{observacao}{Observa{\c c}{\~a}o:}[section]
\newtheorem{observacoes}{Observa{\c c}{\~o}es:}[section]
\newtheorem*{solucao}{Solu{\c c}{\~a}o:}
\newtheorem{proposicao}{Proposi\c{c}\~ao}
\newtheorem{lema}{Lema}
\newtheorem{teorema}{Teorema}
\newtheorem{corolario}{Corol\'ario}
\newenvironment{prova}[1][Prova]{\noindent\textbf{#1:} }{\qedsymbol}%{\ \rule{0.5em}{0.5em}}

\begin{document}

    \begin{figure}[h]
        \begin{minipage}[c]{1.7cm}
            \includegraphics[width=1.7cm]{unb.pdf}
        \end{minipage}
        \hspace{0pt}
        \begin{minipage}[c]{4in}
            {Universidade de Brasília} \\
            {Departamento de Matemática}
        \end{minipage}
    \end{figure}
    \vesp{-1}
    \hrule
    \begin{center}
        {\Large\bf \disciplina\ - Turma \turma}  \\
        \vesp{0.2} {\large\bf Teste \numeroteste\ - Módulo\ \modulo\ -\ \dataavaliacao}
    \end{center}

    \noindent{\bf \underline{Nome: \hspace{9.8cm} Mat.:\qquad \
            \hspace{2.1cm}}}\\
    \vspace*{.01cm}

    \noindent{\bf \underline{Nome: \hspace{9.8cm} Mat.:\qquad \
            \hspace{2.1cm}}}
    \vspace{.4cm}

    \questao{} Seja $S \subseteq M_2(\real)$ dado por:
    \[
    S = \left\{
    \begin{bmatrix}
        a & b\\
        c & d
    \end{bmatrix}
    \mid c = a+b \mbox{ e } d = a
    \right\}.
    \]

    \begin{enumerate}[label={\alph*})]
        \item Mostre que $S$ é um subespaço vetorial de $M_2(\real)$.

        \item Qual a dimensão de $S$?

        \item O conjunto
        \[
        \left\{
        \begin{bmatrix}
            1 & -1\\
            0 & \phantom{x} 1
        \end{bmatrix},
        \begin{bmatrix}
            2 & 1\\
            3 & 4
        \end{bmatrix}
        \right\}
        \]
        é uma base de $S$? Justifique.
    \end{enumerate}

    \solucao Primeiro observe que o conjunto $S$ pode ser escrito como
    \[
        S = \left\{\begin{bmatrix}
            a & b\\a + b & a
        \end{bmatrix} \mid a, b \in \real\right\}.
    \]
    \begin{enumerate}[label={\alph*})]
        \item Tomando $a = b = 0$ segue que
        \[
            \begin{bmatrix}
                0 & 0\\ 0 & 0
            \end{bmatrix} \in S.
        \]
        Logo $S$ é um conjunto não vazio.

        Agora sejam $u_1$, $u_2 \in S$ e $\lambda \in \real$. Temos
        \begin{align*}
            u_1 = \begin{bmatrix}a & b\\a + b & a\end{bmatrix}\\
            u_2 = \begin{bmatrix}c & d\\c + d & c\end{bmatrix}\\
        \end{align*}
        daí:
        \begin{enumerate}[label={\roman*})]
            \item $u_1 + u_2 = \begin{bmatrix}a & b\\a + b & a\end{bmatrix} + \begin{bmatrix}c & d\\c + d & c\end{bmatrix} = \begin{bmatrix}a + c & b + d\\a + b + c + d & a + c\end{bmatrix}\in S$;
            \item $\lambda\cdot u_1 = \lambda\cdot \begin{bmatrix}a & b\\a + b & a\end{bmatrix} = \begin{bmatrix}\lambda a & \lambda b\\\lambda(a + b) & \lambda a\end{bmatrix} = \begin{bmatrix}\lambda a & \lambda b\\\lambda a + \lambda b & \lambda a\end{bmatrix} \in S$.
        \end{enumerate}
        Logo $S$ é de fato um subespaço de $M_2(\real)$.

        \item Para encontrar a dimensão de $S$ precisamos encontrar um base para esse espaço. Para isso, seja $\begin{bmatrix}a & b\\ a + b & a\end{bmatrix} \in S$ um elemento qualquer. Podemos escrever
        \begin{align*}
            \begin{bmatrix}a & b\\a + b & a\end{bmatrix} &= \begin{bmatrix}a & 0\\a & a\end{bmatrix} + \begin{bmatrix}0 & b\\b & 0\end{bmatrix} = a\begin{bmatrix}1 & 0\\1 & 1\end{bmatrix} + b\begin{bmatrix}0 & 1\\1 & 0\end{bmatrix}.
        \end{align*}
        Com isso todo elemento de $S$ é uma combinação linear de $\begin{bmatrix}1 & 0\\1 & 1\end{bmatrix}$ e de $\begin{bmatrix}0 & 1\\1 & 0\end{bmatrix}$. Portanto o conjunto
        \[
            \mathcal{A} = \left\{\begin{bmatrix}1 & 0\\1 & 1\end{bmatrix}, \begin{bmatrix}0 & 1\\1 & 0\end{bmatrix}\right\}
        \]
        gera $S$. Para ser uma base basta que seja L.I. e isso é verdade uma vez que esse conjunto possui somente 2 vetores e é fácil ver que um não é múltiplo escalar do outro. Logo o conjunto $\mathcal{A}$ é uma base para $S$ e como esse conjunto possui 2 vetores então
        \[
            \dim_\real S = 2.
        \]

        \item Como $S$ tem dimensão 2, então qualquer conjunto L.I. contendo 2 vetores será uma base para $S$. Nesse caso é fácil ver que os vetores no conjunto
        \[
        \left\{
        \begin{bmatrix}
            1 & -1\\
            0 & \phantom{x} 1
        \end{bmatrix},
        \begin{bmatrix}
            2 & 1\\
            3 & 4
        \end{bmatrix}
        \right\}
        \]
        são L.I. pois um não é múltiplo escalar no outro, no entanto esse conjunto \textbf{NÃO É UMA BASE} para $S$ pois $ \begin{bmatrix}2 & 1\\3 & 4\end{bmatrix} \notin S$.

        Um conjunto para ser base de um espaço ou subespaço além de ser L.I. e gerar também deve conter vetores desse espaço ou subespaço.
    \end{enumerate}
\end{document}